% ===========================================================================
\chapter{A sense of direction: A*}\label{ch:astar}

Dijkstra's algorithm explores in order of cost so far, so it spreads in
every direction alike. But we know where the target is. A node a kilometre
from the source and two from the target is a better bet than one a
kilometre from the source and five from the target. In 1968 Peter Hart,
Nils Nilsson and Bertram Raphael, building a robot called Shakey at the
Stanford Research Institute, turned that thought into an algorithm, and
called it A*.

\section{The idea}

A* is Dijkstra's algorithm with one change: the priority of a frontier
node $n$ is not just $g(n)$, the cost of the best way to it so far, but
\[
  f(n) = g(n) + h(n),
\]
where $h(n)$ is an \emph{estimate} of the cost still to come, from $n$ to
the target. $f(n)$ is an estimate of the cost of the whole route through
$n$. Nodes on the way to the target have low estimates, so they come out of
the queue early; nodes leading away have high ones, and wait, very often
until the target has been found and the search is over.

That raises a worry. Dijkstra's algorithm is right because it explores
nodes in order of their true cost. If the estimate pushes a node to the
front of the queue, could A* finish the target before finding its cheapest
way? Not if the estimate is never too high.

\begin{mathbox}[When A* finds the best route]
Call an estimate \emph{admissible} if it never overestimates:
$h(n)$ is at most the true cost of the cheapest way from $n$ to the
target. When the target $t$ comes out of the queue, its priority is
$f(t) = g(t) + 0$, the cost of the route found. Every node $m$ still in the
queue has $f(m) \ge f(t)$, and $f(m)$ is \emph{at most} the cost of the
best route through $m$, since $h(m)$ doesn't overestimate. So no route
through any waiting node can be cheaper than the one found.

Our estimates have a slightly stronger property, called being
\emph{consistent}: for every edge from $n$ to $m$, $h(n) \le
\text{cost}(n, m) + h(m)$. The estimate never drops by more than the edge
costs. With a consistent estimate, a node is also finished the first time
it comes out of the queue, exactly as in Dijkstra's algorithm, so the same
code works.
\end{mathbox}

With $h(n) = 0$ for every node, A* is Dijkstra's algorithm. The better the
estimate, the fewer nodes A* explores; a perfect estimate would lead it
straight to the target without a single wasted step.

\section{Estimates}

When the cost is distance, the obvious estimate is the straight-line
distance to the target. No road between two points is shorter than the
straight line between them, so it never overestimates. And it is
consistent, because of the \emph{triangle inequality}: going from $n$
straight to the target is never longer than going from $n$ to $m$ and then
straight to the target.

When the cost is time, the straight-line distance must be turned into a
time that is never too much. Dividing it by the fastest speed on any road
does it: no trip can be quicker than travelling the straight line at top
speed all the way. That is what the graph's \code{top_speed} is for.

Each estimate is a function from a node to a number, built for one target.
Note how the anonymous function uses \code{goal} and \code{speed}, which
are local variables of the function that made it. The anonymous function
\emph{captures} them, and keeps them for as long as it lives, after
\code{time_estimate} itself has returned. A function value with captured
variables is called a \emph{closure}.

\begin{append}{search.nex}
-- A straight-line estimate of the cost from any node to `target`, when
-- the cost is distance: no road between two points is shorter than the
-- straight line between them.
function distance_estimate(g: Road_Graph, target: Integer): Function(Integer): Real do
  let goal: Point := g.island.point(target)
  result := fn (n: Integer): Real do
    result := g.island.point(n).distance_to(goal)
  end
end

-- The same when the cost is time: no trip is quicker than the straight
-- line at the top speed of any road.
function time_estimate(g: Road_Graph, target: Integer): Function(Integer): Real do
  let goal: Point := g.island.point(target)
  let speed: Real := g.top_speed
  result := fn (n: Integer): Real do
    result := g.island.point(n).distance_to(goal) / speed
  end
end
\end{append}

\section{The algorithm}

Here is A*. The only differences from \code{dijkstra} are the extra
parameter, \code{estimate}, and the priority given to each node added to
the queue: \code{c + estimate(e.target)} instead of \code{c}.

\begin{append}{search.nex}
-- A* ("A star"): Dijkstra's algorithm with a sense of direction.
--
-- The priority of a waiting node is the cost of the best way to it so
-- far, plus an estimate of the cost from it to the target. Nodes that
-- lead towards the target come out sooner, and nodes that lead away wait,
-- often until the target is found and they are never explored at all.
-- As long as the estimate is never more than the true cost, the way A*
-- finds is just as cheap as the one Dijkstra finds.
function a_star(g: Road_Graph, source, target: Integer,
  cost: Function(Edge): Real, estimate: Function(Integer): Real): Search_Result
require
  source_on_roads: g.has_node(source)
  target_on_roads: g.has_node(target)
do
  let waiting: Min_Heap[Waiting] := create Min_Heap[Waiting].from_comparator(by_priority())
  waiting.insert(create Waiting.make(source, 0.0, estimate(source)))
  let best: Map[Integer, Real] := {source: 0.0}
  let parents: Map[Integer, Integer] := {}
  let finished: Set[Integer] := #{}
  let explored: Array[Integer] := []
  let found: Boolean := false
  from
  until
    found or waiting.is_empty
  do
    let w: Waiting := waiting.extract_min
    if not finished.contains(w.node) then
      finished.add(w.node)
      explored.add(w.node)
      if w.node = target then
        found := true
      else
        across g.edges_from(w.node) as e do
          let c: Real := w.cost + cost(e)
          if not best.contains_key(e.target) or c < best.get(e.target) then
            best.put(e.target, c)
            parents.put(e.target, w.node)
            waiting.insert(create Waiting.make(e.target, c, c + estimate(e.target)))
          end
        end
      end
    end
  end
  let path: Array[Integer] := []
  if found then
    path := trace_back(parents, source, target)
  end
  result := create Search_Result.make("A*", path,
    best.try_get(target, 0.0), explored, parents)
ensure
  valid_path: not result.found() or is_route(g, result.path)
  cost_adds_up: not result.found()
    or (route_cost(g, result.path, cost) - result.cost).abs < 0.001
end
\end{append}

We could have written \code{dijkstra} as a call to \code{a_star} with an
estimate of zero. Keeping both, side by side, makes it easy to see how
little separates them.

Last, one function that chooses an algorithm by name and plans the
quickest route, or the shortest. The programs at the end of the tutorial
use it, so that you can switch algorithms while they run.

\begin{append}{search.nex}
-- Search with the algorithm called `algorithm` ("bfs", "dijkstra" or
-- "astar"), for the quickest way, or the shortest one if `shortest`.
function find_route(g: Road_Graph, source, target: Integer, algorithm: String,
  shortest: Boolean): Search_Result
require
  source_on_roads: g.has_node(source)
  target_on_roads: g.has_node(target)
  known_algorithm: algorithm = "bfs" or algorithm = "dijkstra" or algorithm = "astar"
do
  let cost: Function(Edge): Real := when shortest then by_distance() else by_time() end
  let estimate: Function(Integer): Real :=
    when shortest then distance_estimate(g, target) else time_estimate(g, target) end
  if algorithm = "bfs" then
    result := breadth_first(g, source, target)
  elseif algorithm = "dijkstra" then
    result := dijkstra(g, source, target, cost)
  else
    result := a_star(g, source, target, cost, estimate)
  end
end
\end{append}

\section{Dijkstra against A*}

\begin{tryit}{try_astar.nex}
intern Try_Pictures

let island: Island_Map := create Island_Map.load(user_path("kestrel_island.map"))
let g: Road_Graph := create Road_Graph.build(island, Mode.Driving)
if ?island.find_place("beach") as a and ?island.find_place("lighthouse") as b then
  let source: Integer := g.nearest_node(island.point(a.node_id))
  let target: Integer := g.nearest_node(island.point(b.node_id))
  console.print_line(summary(g, dijkstra(g, source, target, by_time())))
  let found: Search_Result := a_star(g, source, target, by_time(), time_estimate(g, target))
  console.print_line(summary(g, found))
  picture_of_search(create Map_View.make(island, 1000, 700), found, "astar.png")
end

console.print_line("")
console.print_line("Nodes explored for the quickest route by car:")
across [["harbour", "lighthouse"], ["campsite", "fisher"], ["school", "peak"],
    ["mill", "beach"]] as trip do
  if ?island.find_place(trip.get(0)) as a and ?island.find_place(trip.get(1)) as b then
    let s: Integer := g.nearest_node(island.point(a.node_id))
    let t: Integer := g.nearest_node(island.point(b.node_id))
    console.print_line("  " + (a.name + " to " + b.name).pad_end(" ", 34)
      + dijkstra(g, s, t, by_time()).explored.length.to_string.pad_start(" ", 4)
      + a_star(g, s, t, by_time(), time_estimate(g, t)).explored.length.to_string.pad_start(" ", 5))
  end
end
\end{tryit}
\tryoutput{try_astar}
\mappicture{astar}

A* finds the same route as Dijkstra's algorithm, exactly as quick, but
explores less than half as many nodes: the purple now leans towards the
lighthouse, and most of the town, off to the east of the car park, is
left alone.

The other trips show that the saving varies. From the mill to the beach,
A* explores a fifth as many nodes as Dijkstra's algorithm. But from the
school to Kestrel Peak it saves hardly anything. The quickest way to the
peak ends on the slow tracks through the wood, and the estimate, which
assumes 60\,km/h all the way, is far below the true time. When the
estimate is that far off, A* has little to go on, and behaves almost like
Dijkstra's algorithm.

That is the general rule: the closer the estimate is to the truth, the
less A* explores. Our estimate for time is cautious, because it has to be:
it must never be too high, so it assumes the top speed everywhere. When
walking, every road has the same speed, so the straight-line estimate is
much closer to the truth, and A* is even more focused. The checks in the next chapter compare the two over every pair of
places.
